\section{Introduction}
An Agent receives the instruction \emph{process the next support ticket}. It reads a request, checks an order and approval, consults a model, and submits a refund. Encryption hides messages from a network observer but leaves contacted services, sizes and timing visible. The tool's answer also becomes input to another service: the model. A privacy decoy can cross a second boundary if it invokes a business tool. The requested refund may succeed while an unrelated order or shared quota changes. Figure~\ref{fig:chain} follows these two boundaries.

\begin{figure*}[t]
\centering
\begin{tikzpicture}[x=1cm,y=1cm,font=\small,>=Latex]
\tikzset{box/.style={draw,rounded corners,align=center,minimum height=.85cm,text width=2.4cm},
 note/.style={align=center,font=\footnotesize,text width=3.1cm}}
\node[box,fill=blue!5] (u) at (1.3,1.15) {User\\``Process the next ticket''};
\node[box,fill=blue!5] (a) at (5.1,1.15) {Agent\\plan, check, act, retry};
\node[box,fill=green!6] (m) at (9.2,2.6) {Model service\\next action / streamed output};
\node[box,fill=green!6] (t) at (9.2,-.35) {Tool service\\ticket, order, refund};
\node[box,fill=orange!8] (d) at (14.1,1.15) {Specialist Agent\\optional delegation};
\draw[<->,thick] (u)--node[above,yshift=15pt]{client link C}(a);
\draw[<->,thick] (a.north east)--node[above,sloped]{model link M}(m.west);
\draw[<->,thick] (a.south east)--node[below,sloped]{tool link T}(t.west);
\draw[<->,thick] (a.east)--node[above]{coordination A}(d.west);
\node[note] at (2.0,-.35) {\textbf{Transit observer}\\lengths, gaps, connections;\\destinations when available};
\node[note] at (13.7,-.35) {\textbf{Recipient / active service}\\received arguments;\\optionally selected responses};
\node[draw=black!45,fill=black!3,align=left,text width=15.7cm,inner sep=6pt] at (7.9,-1.65) {
\textbf{Two protection questions:} Where does tool feedback appear next? What work does privacy cover execute before it is discarded? The original task and each observable service remain part of the protection decision.};
\end{tikzpicture}
\caption{Running support-ticket example and observation rights. A tool result returns to the Agent and can enter the next model request. Each encrypted link exposes its own traffic shape; a recipient additionally sees the request addressed to it. Concurrent tickets may share a connection on either link. The arrows represent application communication, not a one-to-one mapping from tool events to packets.}
\label{fig:chain}
\end{figure*}


\textbf{Protection follows the execution across services and state.} Classical traffic analysis supplies observer models, segmentation and adaptive evaluation. We connect twenty-two mechanism anchors to fifteen direct Agent studies and forty-four settings to determine which conditions carry over and which need an application-level constraint. The resulting SoK centers on two findings.

\paragraph{Feedback leaves separately observable representations}
Research queries can be recovered from tool destinations, and task labels can be recognized on model connections~\cite{jeong2026wra,li2026ant,pouryousef2026traffic}. A contacted tool additionally sees its arguments. These observations place protection at different boundaries: tool-side padding leaves later model requests outside its scope, while transit encryption leaves the recipient's arguments visible. Our map connects each inference to the rights and execution stage that make it possible. A paired forwarding pilot measures first-feedback and complete-execution exposure; its adjusted primary comparisons remain inconclusive.

\paragraph{Task completion can conceal executed-cover damage}
Cover discarded at a transport endpoint consumes bandwidth; cover executed by an Agent also changes application work. WRA's five decoys perform additional research and lower query recovery, but its task-utility scores belong to a separate visibility intervention~\cite{jeong2026wra}. This leaves a concrete evaluation question: what happens to the original task and surrounding state on the protected runs?

We test that question with Qwen3:8B and Llama3.1:8B selecting actual tool actions through LangGraph and MCP. Six paired conditions compare no cover, proxy discard, read-only cover, isolated execution, live-state writes and isolated responses visible to the planner. Under the constructed live-write policy, 95/120 Qwen and 73/120 Llama test executions completed the requested task while leaving incorrect global state. Without cover, those counts were 0/120 and 13/120. Both state-integrity reductions survive multiplicity correction; only Qwen's completion reduction does. The result identifies what a completion-only evaluation misses, within a controlled adverse-write design.

Website fingerprinting (WF) and machine-learning security already establish realistic populations, segmentation, transfer and adaptive evaluation~\cite{juarez2014critical,wang2016realistic,dyer2012peekaboo,mathews2023critical,arp2022dos,tramer2020adaptive}. Our contributions are their \emph{transfer and deployment maps} for Agent execution, a \emph{source-bound account of remaining views and executed cover}, and a \emph{matched measurement of inference, task completion and state integrity}. Among four published or preprint query/workflow privacy settings from two sources, adaptation is located in two and original-task utility in one; none documents both. The new controlled study supplies all three endpoints for its own workload and keeps its results separate from this literature count.
